%%%%%%%%%%%%%%%%%%%%%%%%%%%%%%%%%%%%%%%%%%%%%%%%%%%%%%%%%%%%%%%%%%%%%%%%%%%%%%%%
\chapter{Conclusion}\label{ch:conclusion}
%%%%%%%%%%%%%%%%%%%%%%%%%%%%%%%%%%%%%%%%%%%%%%%%%%%%%%%%%%%%%%%%%%%%%%%%%%%%%%%%
In the established literature on Curry-Howard interpretations for linear logic and
$\pi$-calculi, elimination rules have seen little use, with the introduction
rules of Girard's linear logic being the standard foundation.
With $\piFD$, we have presented not only the operational semantics of a calculus
that relies solely on elimination rules, but also a new approach to the design
of languages for concurrent processes. By re-examining the relationship between
axioms, elimination rules, and the rule of \textsc{cut}, we have established
a novel computational interpretation, one where the structural rule
of \textsc{cut} serves to compose concurrent processes, but is disentangled
from the actual act of communication between channel endpoints.
With a fully formalized metatheory in Rocq,
we have thus provided a rigorous alternative, expanding the landscape
of session-typed process calculi.

\section{Related and Future Work}
\paragraph{Computational sequent calculi.}
Over the years, multiple calculi have emerged in the context of
Curry-Howard isomorphisms for classical logic, such as
Parigot's $\lambda\mu$-calculus \cite{Parigot1992:lambda-mu} or
Curien and Herbelin's $\bar{\lambda}\mu\tilde{\mu}$-calculus \cite{Curien2000:DualityOfComputation}.
Naturally, the question arises as to how these calculi relate to
linear interpretations of classical calculi. A relation between the
linear $\lambda\mu$-calculus and Wadler's $\CP$ was developed by Paykin \cite{Paykin2016:lambda-mu-is-cp}.
A similar development can be carried out for $\piFD$ and a linear $\lambda\mu\tilde{\mu}$-calculus, where we
relate parallel compositions with critical pairs, and links with cuts between
producers and consumers that are neither $\mu$- nor $\tilde{\mu}$-abstractions.
We leave the formal details of this translation to future work.

\paragraph{Verified Session-Typed Communication.}
Session types allow programmers to structure communication by specifying
type-level protocols. To make these protocols more expressive and enable
the verification of different properties within concurrent programs,
a line of work has studied dependent session types
\cite{Pfennig2011:ProofCarryingCodeInSessionCalculi, Toninho2011:DependentSessionTypesIntuitionisticTypeTheory}.
Other approaches leveraging index refinements to provide safety guarantees
have also been developed, such as in the language Granule \cite{Orchard2019:granule}.

\paragraph{Non-Determinism.}
One of the main results of this work is the mechanized proof that $\piFD$
enjoys the Church-Rosser property, ensuring deterministic evaluation.
However, modeling concurrency often requires reasoning about races and non-determinism.
Principled approaches to extending linear logic with non-deterministic behavior
have been explored in the literature, including the introduction of client pools
by \citet{Kokke2019:RacesInLinearLogic} and coexponentials by \citet{Qian2021:CoExponentials}.

\paragraph{Multiparty Session Types.}
Binary session types, introduced by \citet{Honda1993:TypesForDyadicInteraction}, form the basis of communication in $\piFD$.
Modern concurrent systems often involve protocols among more than two participants.
This has led to the development of Multiparty Session Types (MPST)
introduced by \citet{Honda2008:MPST}, which remains an active area of research.
